% Source: https://artofproblemsolving.com/community/c3h1844338_parabolas
% Stable identity: aops-topic-1844338-q01
\textbf{(a)} Write $y=ax^2+bx+c$. The point $O$ gives $c=0$, and $P$, $Q$ give $a-b=\sqrt3$ and $a+b=\sqrt3$, so $b=0$ and $a=\sqrt3$. The parabola is $\boxed{y=\sqrt3\,x^2}$, with vertex $O$ and a vertical axis.

\textbf{(b)} The centre of the triangle is $G=\left(0,\dfrac{2\sqrt3}{3}\right)$. A rotation about $G$ through $120^\circ$ or $240^\circ$ maps the equilateral triangle onto itself, permuting its vertices. A rotation maps a parabola onto a congruent parabola, and it maps the points $O,P,Q$ of $y=\sqrt3x^2$ onto the points of the image, which are again $O,P,Q$ in some order. So the two images are parabolas through $O$, $P$ and $Q$. Their axes are the vertical axis turned through $120^\circ$ and $240^\circ$, and their vertices are the images of $O$, namely $P$ and $Q$. Three parabolas with three different axis directions are three different curves.

\begin{asy}
usepackage("amsmath");
size(6.5cm);
defaultpen(fontsize(10pt));
import graph;
real s3=sqrt(3);
pair O=(0,0), P=(-1,s3), Q=(1,s3), G=(0,2*s3/3);
real f(real x){return s3*x^2;}
path par=graph(f,-1.18,1.18);
draw((-2.6,0)--(2.6,0),gray(0.7),Arrow(5));
draw((0,-1.3)--(0,3.4),gray(0.7),Arrow(5));
draw(O--P--Q--cycle,gray(0.45)+dashed);
draw(par,blue+linewidth(0.9));
draw(rotate(120,G)*par,red+linewidth(0.9));
draw(rotate(240,G)*par,heavygreen+linewidth(0.9));
dot(O);dot(P);dot(Q);dot(G,gray(0.3));
label("$O$",(0,-0.33));label("$P$",P+0.3*dir(232));label("$Q$",Q+0.3*dir(-50));label("$G$",G,E,gray(0.3));
label("$y=\sqrt{3}\,x^2$",f(1.18)*N+1.18*E,E,blue);
label("\mbox{}",(min(currentpicture,user=true)+max(currentpicture,user=true))/2,black); // keeps the text after the figure black when LaTeX typesets the labels inline
\end{asy}

\textbf{(c)} Let the three points be $X_1,X_2,X_3$, not on one line. Choose any line direction $e$ that is not parallel to any of the three sides $X_1X_2$, $X_2X_3$, $X_3X_1$; all but three directions qualify. Rotate the coordinate system so that the new vertical axis (the $t$-axis) is parallel to $e$, and call the new coordinates $(s,t)$. Two of the points have equal $s$-coordinates only if the line through them is parallel to $e$, so the $s$-coordinates $s_1,s_2,s_3$ are different. Then
\[
t=t_1\dfrac{(s-s_2)(s-s_3)}{(s_1-s_2)(s_1-s_3)}+t_2\dfrac{(s-s_1)(s-s_3)}{(s_2-s_1)(s_2-s_3)}+t_3\dfrac{(s-s_1)(s-s_2)}{(s_3-s_1)(s_3-s_2)}
\]
is a quadratic function $t=\alpha s^2+\beta s+\gamma$ whose graph passes through the three points (at $s=s_i$ all fractions vanish except the $i$th, which equals $1$). Its leading coefficient $\alpha$ is not $0$: otherwise the three points would lie on the line $t=\beta s+\gamma$. So the graph is a parabola whose axis is parallel to $e$.

As $s\to\pm\infty$ the points of this parabola go off to infinity in the direction of $e$, so parabolas obtained from different directions $e$ are different curves. There are infinitely many admissible directions, hence infinitely many parabolas through $X_1,X_2,X_3$. (So ``three points determine a parabola'' is true only once the direction of the axis is fixed, for example for parabolas $y=ax^2+bx+c$.)
